% Bedingte Wahrscheinlichkeit und Bayes – bank/bedingte-wahrscheinlichkeit-und-bayes/ (zusammenbau v0.4, Heft)
\einheitenkopf[t2]{Bedingte Wahrscheinlichkeit und Bayes}
\setcounter{aufgabe}{10}

% Anlauf (ohne Prüfkennung)
\begin{aufgabe}{Der Quotient – Anlauf}
\begin{teile}
\teil Gesucht ist der Anteil der Mädchen unter den Vereinsmitgliedern, die Fußball spielen. Durch wessen Anteil wird geteilt? \\ \kreuz{alle Mitglieder} \\ \kreuz{die Fußball spielenden Mitglieder} \\ \kreuz{die Mädchen}
\teil In einem Jahrgang lernen 40\,\% der Schüler Spanisch; 16\,\% aller Schüler lernen Spanisch und spielen ein Instrument. Ein Schüler, der Spanisch lernt, wird zufällig ausgewählt. Mit welcher Wahrscheinlichkeit spielt er ein Instrument? \leerfeld
\teil Befragt wurden 200 Personen. A: „trinkt morgens Kaffee“, B: „steht vor 6 Uhr auf“. Die Vierfeldertafel zeigt die Anzahlen. Eine Person, die vor 6 Uhr aufsteht, wird zufällig ausgewählt. Mit welcher Wahrscheinlichkeit trinkt sie morgens Kaffee?

\vierfeldertafel{A}{B}{36,54,90,24,86,110,60,140,200}
\leerfeld
\end{teile}
\end{aufgabe}

\begin{aufgabe}{Der Quotient – Prüfungsaufgaben}
\begin{teile}
\teil Von den Beschäftigten eines Verkehrsbetriebs arbeiten 35\,\% im Fahrdienst; 14\,\% der Beschäftigten sind Frauen und arbeiten im Fahrdienst; insgesamt sind 40\,\% der Beschäftigten Frauen. Eine zufällig ausgewählte beschäftigte Person ist eine Frau. Mit welcher Wahrscheinlichkeit arbeitet sie im Fahrdienst? \hfill (Abitur 2023 GK) \\ \leerfeld
\teil In einer Klinik sind 81\,\% der Patienten mit dem Essen zufrieden; 4\,\% aller Patienten liegen auf der Kinderstation und sind nicht zufrieden; 10\,\% aller Patienten liegen auf der Kinderstation. Untersuche, ob der Anteil der Unzufriedenen auf der Kinderstation größer ist als auf den übrigen Stationen. \hfill (Abitur 2021 GK)
\end{teile}
\end{aufgabe}

% Anlauf (ohne Prüfkennung)
\begin{aufgabe}{Bayes – Anlauf}
\begin{teile}
\teil Ein Baum hat als erste Stufe das Werk (1 oder 2), als zweite Stufe „defekt“ oder „nicht defekt“. Gefragt ist: Wie wahrscheinlich ist ein Gerät aus Werk 1 defekt? Geht das mit dem Baum oder gegen den Baum? \\ \kreuz{mit dem Baum} \\ \kreuz{gegen den Baum}
\teil Das Baumdiagramm beschreibt eine Untersuchung mit einem Test. Eine Person hat ein positives Ergebnis. Mit welcher Wahrscheinlichkeit ist sie krank?

\baumzwei{krank/0.1,gesund/0.9}{positiv/0.8,negativ/0.2}{positiv/0.1,negativ/0.9}
\leerfeld
\teil In einem Lager stammen 30\,\% der Geräte aus Werk 1, der Rest aus Werk 2; aus Werk 1 sind 4\,\% defekt, aus Werk 2 sind es 2\,\%. Gib an, welche Wahrscheinlichkeit der Term $\frac{0{,}3 \cdot 0{,}04}{0{,}3 \cdot 0{,}04 + 0{,}7 \cdot 0{,}02}$ beschreibt.
\end{teile}
\end{aufgabe}

\begin{aufgabe}{Bayes – Prüfungsaufgaben}
\begin{teile}
\teil In einem Sportverein sind zwei Fünftel der Mitglieder Jugendliche. 6\,\% der Jugendlichen und 9\,\% der Erwachsenen sind mit den Trainingszeiten unzufrieden. Ein zufällig ausgewähltes Mitglied ist unzufrieden. Mit welcher Wahrscheinlichkeit ist es erwachsen? \hfill (Abitur 2020 GK) \\ \leerfeld
\teil In einem Tierbestand liegt bei 2\,\% der Tiere eine Infektion vor. Ein Schnelltest ist bei infizierten Tieren mit 95\,\% positiv, bei nicht infizierten mit 3\,\%. Deute den Term $\frac{0{,}02 \cdot 0{,}95}{0{,}02 \cdot 0{,}95 + 0{,}98 \cdot 0{,}03}$ im Sachzusammenhang. \hfill (Abitur 2021 GK)
\teil Eine Druckerei druckt 30\,\% ihrer Aufträge auf Maschine A. Bei 6\,\% aller Aufträge wird auf A gedruckt und es gibt eine Reklamation, bei 9\,\% aller Aufträge wird nicht auf A gedruckt und es gibt eine Reklamation. Ein Auftrag wird reklamiert. Untersuche, ob die Wahrscheinlichkeit, dass er nicht auf A gedruckt wurde, größer als 50\,\% ist. \hfill (Abitur 2026 GK)
\teil Bei einer Befragung in einem Verein erhalten 90\,\% der Mitglieder verdeckt die Frage „Willst du den Verein verlassen?“, die übrigen die Frage „Willst du im Verein bleiben?“; jede Person antwortet ehrlich. 15\,\% aller Mitglieder wollen den Verein verlassen. Eine Person mit der Antwort Ja wird zufällig ausgewählt. Mit welcher Wahrscheinlichkeit will sie den Verein verlassen? \hfill (Abitur 2021 GK) \\ \leerfeld
\end{teile}
\end{aufgabe}

% Anlauf (ohne Prüfkennung)
\begin{aufgabe}{Mit Parameter – Anlauf}
\begin{teile}
\teil Wo steht der Parameter im Term $\frac{0{,}3}{0{,}3 + a}$? \\ \kreuz{nur im Zähler} \\ \kreuz{nur im Nenner} \\ \kreuz{in Zähler und Nenner}
\teil Ein Anteil a der Kunden eines Geschäfts, $0 < a < 1$, sind Stammkunden. 60\,\% der Stammkunden und 10\,\% der übrigen Kunden kaufen ein Angebot. Ein Kunde kauft das Angebot. Stelle einen Term in a für die Wahrscheinlichkeit auf, dass er Stammkunde ist. \leerfeld
\teil Ein Drittel der Kunden eines Versandhändlers sind Neukunden. 10\,\% der Neukunden und ein Anteil b der Stammkunden reklamieren. Durch bessere Beratung sinkt b. Begründe ohne Rechnung, dass dabei die Wahrscheinlichkeit wächst, dass ein reklamierender Kunde ein Neukunde ist.
\end{teile}
\end{aufgabe}

\begin{aufgabe}{Mit Parameter – Prüfungsaufgaben}
\begin{teile}
\teil Ein Onlinehändler verschickt 70\,\% seiner Pakete mit dem Dienst D, davon kommen 90\,\% pünktlich an; von den übrigen Paketen kommt ein Anteil a nicht pünktlich an. Ein Paket kam nicht pünktlich an. Welcher der Graphen A, B, C zeigt die Wahrscheinlichkeit, dass es mit D verschickt wurde, in Abhängigkeit von a? Begründe ohne Rechnung. \hfill (Abitur 2020 GK)

\begin{ksys}[xmin=0,xmax=1,ymin=0,ymax=1,xstep=0.2,ystep=0.2,ablesen] \funktion{1-\x}{A} \funktion{0.07/(0.07+0.3*\x)}{B} \funktion{0.07/(0.07+0.3*(1-\x))}{C} \end{ksys}
\teil In einer Stadt nutzen 50\,\% der Haushalte Fernwärme; von ihnen haben 30\,\% einen hohen Energieverbrauch, von den übrigen 10\,\%. Durch Sanierungen sinkt der Anteil der Fernwärme-Haushalte mit hohem Verbrauch um x. Die Funktion $f(x) = \frac{0{,}5 \cdot (0{,}3 - x)}{0{,}5 \cdot (0{,}3 - x) + 0{,}05}$ mit $0 \le x \le 0{,}3$ gibt die Wahrscheinlichkeit an, dass ein Haushalt mit hohem Verbrauch Fernwärme nutzt. Bestimme am Graphen x mit $f(x) = 0{,}6$ und deute x und $f(x)$ im Sachzusammenhang. \hfill (Abitur 2025 GK)

\begin{ksys}[xmin=0,xmax=0.3,ymin=0,ymax=1,xstep=0.05,ystep=0.1,ablesen] \funktion{0.5*(0.3-\x)/(0.5*(0.3-\x)+0.05)}{f} \end{ksys}
\end{teile}
\end{aufgabe}
